\scanpage{129}

sous-module \inserted{\unsure{ess. const.|ens. const.|en. const.}} d'un $R^m g_*(\mathbb Q_\ell)$, $g:V\unsure{\text{---}|\longrightarrow}U$ un morphisme \uline{projectif} de dim \inserted{relative} $\leq d$. On a donc des suites exactes

\[
\begin{aligned}
0&\longrightarrow \underset{\diagup}{G}\mathrel{\text{---}} R^{\unsure{m|n}}g_*(\mathbb Q_\ell)\mathrel{\text{---}}H\mathrel{\text{---}}0\\
0&\mathrel{\text{---}}\underset{*}{I}\mathrel{\text{---}}\underset{*}{G}\mathrel{\text{---}}F\longrightarrow 0
\end{aligned}
\]

et \deleted{\unsure{une}} \inserted{de} nouvelles applications du lemme 2 nous \unsure{ramène|ramènent} au cas de $I$ et $G$, $\unsure{\mathbb R|R}^{\unsure{-|n}}f_*(F)$ lorsque \deleted{$I,G$} $F$ ess. const. $\subset R^{m}g_*(\mathbb Q_\ell)$. Pour continuer, il faut utiliser le ``th.'' \uline{de complète réductibilité} pour nous \unsure{ramener|soumis} au cas où $F$ est de la forme $R^m g_*(\mathbb Q_\ell)$

on doit donc examiner le cas

\[ R^p_! f_*(R^q g_*(\mathbb Q_\ell))\qquad p,q\geq0 \]

\deleted{\unsure{1\textsuperscript{o})|b)} Cas $m=0$. Nous procédons par réc. \uline{descendante} sur $p+q$ \deleted{\unsure{suivant}}, en supposant que c'est vrai \unsure{pr|pour} $(p',q')$ si $p'>p$ \uuline{ou} $q$}

% archive sheet number: 128
